%Academic CV LaTeX Template
% Author: Dubasi Pavan Kumar
% LinkedIn: https://www.linkedin.com/in/im-pavankumar/
% License: MIT
%
% For errors, suggestions, or improvements, please contact:
% Email: pavankumard.pg19.ma@nitp.ac.in

\documentclass[a4paper,11pt]{article}

% Package imports
\usepackage{latexsym}
\usepackage{xcolor}
\usepackage{float}
\usepackage{ragged2e}
\usepackage[empty]{fullpage}
\usepackage{wrapfig}
\usepackage{lipsum}
\usepackage{tabularx}
\usepackage{titlesec}
\usepackage{geometry}
\usepackage{marvosym}
\usepackage{verbatim}
\usepackage{enumitem}
\usepackage{fancyhdr}
\usepackage{multicol}
\usepackage{graphicx}
\usepackage{cfr-lm}
\usepackage[T1]{fontenc}
\usepackage{fontawesome5}

% Color definitions
\definecolor{darkblue}{RGB}{0,0,139}

% Page layout
\setlength{\multicolsep}{0pt}
\pagestyle{fancy}
\fancyhf{} % clear all header and footer fields
\fancyfoot{}
\renewcommand{\headrulewidth}{0pt}
\renewcommand{\footrulewidth}{0pt}
\geometry{left=1.4cm, top=0.8cm, right=1.2cm, bottom=1cm}
\setlength{\footskip}{5pt} % Addressing fancyhdr warning

% Hyperlink setup (moved after fancyhdr to address warning)
\usepackage[hidelinks]{hyperref}
\hypersetup{
    colorlinks=true,
    linkcolor=darkblue,
    filecolor=darkblue,
    urlcolor=darkblue,
}

% Custom box settings
\usepackage[most]{tcolorbox}
\tcbset{
    frame code={},
    center title,
    left=0pt,
    right=0pt,
    top=0pt,
    bottom=0pt,
    colback=gray!20,
    colframe=white,
    width=\dimexpr\textwidth\relax,
    enlarge left by=-2mm,
    boxsep=4pt,
    arc=0pt,outer arc=0pt,
}

% URL style
\urlstyle{same}

% Text alignment
\raggedright
\setlength{\tabcolsep}{0in}

% Section formatting
\titleformat{\section}{
  \vspace{-4pt}\scshape\raggedright\large
}{}{0em}{}[\color{black}\titlerule \vspace{-7pt}]

% Custom commands
\newcommand{\resumeItem}[2]{
  \item{
    \textbf{#1}{\hspace{0.5mm}#2 \vspace{-0.5mm}}
  }
}

\newcommand{\resumePOR}[3]{
\vspace{0.5mm}\item
    \begin{tabular*}{0.97\textwidth}[t]{l@{\extracolsep{\fill}}r}
        \textbf{#1}\hspace{0.3mm}#2 & \textit{\small{#3}}
    \end{tabular*}
    \vspace{-2mm}
}

\newcommand{\resumeSubheading}[4]{
\vspace{0.5mm}\item
    \begin{tabular*}{0.98\textwidth}[t]{l@{\extracolsep{\fill}}r}
        \textbf{#1} & \textit{\footnotesize{#4}} \\
        \textit{\footnotesize{#3}} &  \footnotesize{#2}\\
    \end{tabular*}
    \vspace{-2.4mm}
}

\newcommand{\resumeProject}[4]{
\vspace{0.5mm}\item
    \begin{tabular*}{0.98\textwidth}[t]{l@{\extracolsep{\fill}}r}
        \textbf{#1} & \textit{\footnotesize{#3}} \\
        \footnotesize{\textit{#2}} & \footnotesize{#4}
    \end{tabular*}
    \vspace{-2.4mm}
}

\newcommand{\resumeSubItem}[2]{\resumeItem{#1}{#2}\vspace{-4pt}}

\renewcommand{\labelitemi}{$\vcenter{\hbox{\tiny$\bullet$}}$}
\renewcommand{\labelitemii}{$\vcenter{\hbox{\tiny$\circ$}}$}

\newcommand{\resumeSubHeadingListStart}{\begin{itemize}[leftmargin=*,labelsep=1mm]}
\newcommand{\resumeHeadingSkillStart}{\begin{itemize}[leftmargin=*,itemsep=1.7mm, rightmargin=2ex]}
\newcommand{\resumeItemListStart}{\begin{itemize}[leftmargin=*,labelsep=1mm,itemsep=0.5mm]}

\newcommand{\resumeSubHeadingListEnd}{\end{itemize}\vspace{2mm}}
\newcommand{\resumeHeadingSkillEnd}{\end{itemize}\vspace{-2mm}}
\newcommand{\resumeItemListEnd}{\end{itemize}\vspace{-2mm}}
\newcommand{\cvsection}[1]{%
\vspace{2mm}
\begin{tcolorbox}
    \textbf{\large #1}
\end{tcolorbox}
    \vspace{-4mm}
}

\newcolumntype{L}{>{\raggedright\arraybackslash}X}%
\newcolumntype{R}{>{\raggedleft\arraybackslash}X}%
\newcolumntype{C}{>{\centering\arraybackslash}X}%

% Commands for icon sizing and positioning
\newcommand{\socialicon}[1]{\raisebox{-0.05em}{\resizebox{!}{1em}{#1}}}
\newcommand{\ieeeicon}[1]{\raisebox{-0.3em}{\resizebox{!}{1.3em}{#1}}}

% Font options
\newcommand{\headerfonti}{\fontfamily{phv}\selectfont} % Helvetica-like (similar to Arial/Calibri)
\newcommand{\headerfontii}{\fontfamily{ptm}\selectfont} % Times-like (similar to Times New Roman)
\newcommand{\headerfontiii}{\fontfamily{ppl}\selectfont} % Palatino (elegant serif)
\newcommand{\headerfontiv}{\fontfamily{pbk}\selectfont} % Bookman (readable serif)
\newcommand{\headerfontv}{\fontfamily{pag}\selectfont} % Avant Garde-like (similar to Trebuchet MS)
\newcommand{\headerfontvi}{\fontfamily{cmss}\selectfont} % Computer Modern Sans Serif
\newcommand{\headerfontvii}{\fontfamily{qhv}\selectfont} % Quasi-Helvetica (another Arial/Calibri alternative)
\newcommand{\headerfontviii}{\fontfamily{qpl}\selectfont} % Quasi-Palatino (another elegant serif option)
\newcommand{\headerfontix}{\fontfamily{qtm}\selectfont} % Quasi-Times (another Times New Roman alternative)
\newcommand{\headerfontx}{\fontfamily{bch}\selectfont} % Charter (clean serif font)

\begin{document}
\headerfontiii

% Header
\begin{center}
    {\Huge\textbf{Amani BENZARTI}}
\end{center}
\vspace{-6mm}

\begin{center}
    \small{
    +1-263-382-9492 | \href{mailto:amani.benzartii@gmail.com}{amani.benzartii@gmail.com}
    }
\end{center}
\vspace{-6mm}

\begin{center}
    \small{8375-4, 24eme Avenue, Montréal, Quebec - H1Z 3Z2, Canada}
\end{center}

%\vspace{-4mm}

\section{\textbf{Profil}}
\vspace{1mm}
\small{
Médecin de famille récemment diplômée, rigoureuse et engagée, ayant acquis une expérience clinique diversifiée en soins de première ligne, urgences et différentes spécialités médicales. Dotée d’excellentes capacités de communication, d’adaptation et de collaboration interprofessionnelle, je place l’écoute, l’empathie et la continuité des soins au cœur de ma pratique. Mon parcours clinique, universitaire et communautaire témoigne de mon intérêt pour une médecine familiale globale, humaine et centrée sur le patient. Bilingue français-anglais, je souhaite poursuivre ma formation en résidence au Canada et mettre mes compétences, ma capacité d’adaptation et mon engagement au service des patients et de leur communauté.
}
%\vspace{-2mm}
%%%%%%%%%%%%%%%%%%%%%%%%%%%%%%%%%%%%%%%%%%%%%%%%%%%%%%%%%%%%%%%%%%%%%%%%%%
\section{\textbf{FORMATIONS ACADÉMIQUES}}
%\vspace{-0.4mm}
\resumeSubHeadingListStart
\resumeSubheading
{Equivalence MD du Collège des Médecins du Québec}{}
{}{ Juillet 2026}
%\vspace{-3mm}
\resumeSubheading
{Titre de Licenciée du Conseil médical du Canada (LCMC)}{}
{}{ Juillet 2026}
%\vspace{-3mm}
\resumeSubheading
{Examens d’équivalence du Conseil Médical du Canada}{}
{}{}
%\vspace{-3mm}
\resumeItemListStart
\item Examen d'aptitude du Conseil médical du Canada (EACMC) : Réussite — Avril 2026
\item Examen de la Collaboration nationale en matière d’évaluation CNE : Réussite — Sep 2025
\resumeItemListEnd

\resumeSubheading
{DIPLOME DE DOCTEUR EN MEDECINE (M.D) }{}
{Faculté de médecine de Tunis, Tunisie}{Mars 2024}
\resumeItemListStart
\item Obtenu avec mention très honorable et félicitations du jury
\resumeItemListEnd

\resumeSubheading
{CERTIFICAT D'ÉTUDES COMPLEMENTAIRES EN THÉRAPIES BRÈVES}{}
{Faculté de médecine de Tunis, Tunisie}{2023 - 2025}
\resumeItemListStart
\item Obtenu avec mention très honorable et félicitations du jury
\resumeItemListEnd
%\vspace{1mm}
\resumeSubheading
{RÉUSSITE À L’EXAMEN CLASSANT NATIONAL (ECN) }{}
{Faculté de médecine de Tunis, Tunisie}{2021}
%\vspace{1mm}
%\resumeSubheading
%{DIPLOME D’ETUDES COLLEGIALES: SECTION MATHÉMATIQUES }{}
%{Lycée pilote ARIANA (LPA) }{2011}
\resumeSubHeadingListEnd
%\vspace{-4mm}
%%%%%%%%%%%%%%%%%%%%%%%%%%%%%%%%%%%%%%%%%%%%%%%%%%%%%%%%%%%%%%%%%%%%%%%%%%
\section{\textbf{EXPÉRIENCES PROFESSIONNELLES}}
%\vspace{-0.4mm}
  \resumeSubHeadingListStart
    \resumeSubheading
      {{STAGIAIRE OBSERVATRICE – SERVICE O.R.L} - CHUM}{}
      {}{Septembre 2026}
      \vspace{-3mm}
      \resumeItemListStart
        \item Observation de l’évaluation, du diagnostic et de la prise en charge des pathologies ORL
        \item Reconnaissance des signes d’alarme et les indications de référence en ORL, pertinentes \newline à la pratique de la médecine
    \resumeItemListEnd
    %\resumeSubheading
     % {{STAGIAIRE OBSERVATRICE – Médecine familiale - Clinique Médicale Privée Mont Royal}}{}
     % { }{Août 2026}
    %  \vspace{-3mm}
   %   \resumeItemListStart
        %\item Familiarisation avec la prévention, le dépistage et la prise en charge des maladies aiguës\newline et chroniques
    %    \item Observation de l’approche centrée sur le patient, de la collaboration interprofessionnelle \newline et du fonctionnement des soins de première ligne au Canada.
    %\resumeItemListEnd

    \resumeSubheading
      {{STAGIAIRE OBSERVATRICE – SERVICE DES URGENCES - Hôpital Maisonneuve-Rosemont}}{}
      {}{Avril 2025}
      \resumeItemListStart
        \item Observation de l’évaluation et de la prise en charge de patients présentant diverses conditions \newline aiguës.
        \item Familiarisation avec le fonctionnement d’un service d’urgence au Québec.
        %\item Observation de la collaboration interprofessionnelle et de la communication médecin-patient.
        \item Familiarisation avec le dossier médical informatisé et la documentation clinique.
    \resumeItemListEnd

    \resumeSubheading
      {{DOCTORANTE - Faculté de médecine de Tunis}}{}
      {}{2024}
      \resumeItemListStart
        \item Étude portant sur la valeur pronostique des lymphocytes infiltrant les tumeurs dans le \newline cancer gastrique.
        \item Article publié en 2026 sous le titre de “Prognostic value of tumor-infiltrating lymphocytes \newline and their association with clinical and pathological factors in gastric cancer” (ABCD Arq \newline Bras Cir Dig. 2026;39. PMID: 42585393. DOI: 10.1590/0102-672020260000019e1948.)
        \item Communication Orale présentée sous le nom « Les lymphocytes infiltrant la tumeur (TILS):\newline valeur pronostic et association avec les facteurs cliniques et pathologiques dans le cancer \newline gastrique » lors du congrès annuel de l'Association Française de Chirurgie (AFC), en \newline Septembre 2024, Paris, France.
    \resumeItemListEnd
    \vspace{5mm}
    \hspace{-2.5mm}{\textbf{• RÉSIDENTE EN MÉDECINE FAMILIALE} - \textit{Ministère de la santé publique de Tunis}}
    \vspace{1mm}
    \resumeSubheading
      {{Médecine familiale}}{}
      {Centre de santé de base La Marsa}{Janvier - Décembre 2024}
      \resumeItemListStart
        \item Consultations de médecine familiale et suivi longitudinal des patients de tout âge.
        \item Prise en charge des maladies aiguës courantes et suivi des maladies chroniques.
        \item Santé de la femme: planification familiale, ITS, dépistage des cancers du col et du sein.
        \item Santé de l’enfant: consultations pédiatriques, vaccination et promotion de l’allaitement \newline maternel.
        \item Prévention, dépistage et éducation à la santé.
    \resumeItemListEnd

    \resumeSubheading
      {{Endocrinologie}}{}
      {Institut National Zouhaier Kallel de Nutrition}{Octobre - Décembre 2023}
      \resumeItemListStart
        \item Prise en charge et suivi des patients atteints de diabète et autres pathologies endocriniennes \newline et métaboliques.
        \item Dépistage et suivi des complications aiguës et chroniques du diabète.
        %\item Évaluation nutritionnelle et prise en charge de l’obésité.
    \resumeItemListEnd

    \resumeSubheading
      {{Pneumologie}}{}
      {Hôpital Mami Ariana }{Juillet - Septembre 2023}
      \resumeItemListStart
        \item Prise en charge et suivi de patients présentant des pathologies respiratoires aiguës \newline et chroniques.
        \item Interprétation des examens complémentaires respiratoires et participation à l’élaboration \newline des stratégies diagnostiques et thérapeutiques.
    \resumeItemListEnd

    \resumeSubheading
      {{Médecine de travail}}{}
      {Hôpital La Rabta}{Avril - Juin 2023}
      \resumeItemListStart
        \item Évaluation de l’aptitude au travail, suivi de la santé des travailleurs et dépistage des \newline pathologies liées aux expositions professionnelles.
    \resumeItemListEnd

    \resumeSubheading
      {{Médecine préventive}}{}
      {Observatoire National des Maladies Nouvelles et émergentes}{Janvier - Mars 2023}
      \resumeItemListStart
        \item Initiation à la surveillance épidémiologique, à la veille sanitaire et à l’investigation des \newline maladies émergentes et des risques pour la santé publique.
        \item Participation active en tant qu’enquêtrice à l’évaluation post-introduction du vaccin \newline contre la COVID-19, menée en collaboration avec l’Organisation mondiale de la Santé (OMS).
        \item Contribution à la veille sanitaire internationale via la plateforme EIOS : surveillance, \newline collecte et analyse d’informations relatives aux événements sanitaires émergents.
    \resumeItemListEnd

    \resumeSubheading
      {{Psychiatrie }}{}
      {Hôpital Razi}{Octobre - Décembre 2022}
      \resumeItemListStart
        \item Participation à l’évaluation diagnostique et à la prise en charge globale des patients \newline présentant divers troubles psychiatriques.
        \item Surveillance de la réponse thérapeutique, de la tolérance et des effets indésirables \newline des traitements psychotropes.
    \resumeItemListEnd
    \vspace{-1mm}
    \resumeSubheading
      {{Dermatologie}}{}
      {Hôpital militaire}{Août - Septembre 2022}
    \resumeSubheading
      {{ORL}}{}
     {Hôpital Monji Slim}{Juillet - Aout 2022}
    \resumeSubheading
      {{Urgences }}{}
      {Hôpital regional de Ben Arous}{Janvier - Juin 2022}
      \resumeItemListStart
        \item Service des urgences : triage, unité d’hospitalisation de courte durée et soins intensifs.
        \item Évaluation au triage et priorisation des patients selon la gravité clinique.
        \item Participation à la prise en charge initiale des urgences médicales.
        \item Prise en charge des cas suspects de COVID-19 au sein de l’unité de dépistage dédiée.
        \item Participation à la surveillance et au suivi des patients nécessitant une hospitalisation \newline de courte durée.
        \item Aide à la prise en charge des patients critiques en unité de soins intensifs.
    \resumeItemListEnd
    \vspace{4mm}
    \hspace{-1.5mm}{\textbf{\underline{INTERNE EN MEDECINE}} *}
    \vspace{2mm}

    \resumeSubheading
      {{Chirurgie Générale}}{Octobre - Décembre 2019}
      {Hôpital régional de Nabeul}{}

    \resumeSubheading
      {{Cardiologie}}{Juillet - Septembre 2019}
      {Hôpital Charles Nicole de Tunis}{}

    \resumeSubheading
      {{Gynécologie obstétrique}}{Avril - Juin 2019}
      {Hôpital Régional de Nabeul}{}

    \resumeSubheading
      {{Pédiatrie }}{Janvier - Mars 2019}
      {Hôpital Régional de Nabeul}{}

  \resumeSubHeadingListEnd

  \hline

  \vspace{3mm}
  \scriptsize{* Internat rotatoire d’un an avec stages cliniques supervisés, permettant l’acquisition d’une expérience multidisciplinaire et de compétences essentielles à une prise en charge globale.}

\vspace{-2mm}

\section{\textbf{PARTICIPATIONS AUX CONGRÈS SCIENTIFIQUES} }
\vspace{0.2mm}
\small{
\begin{itemize}
\item 28eme Congrès de Pneumologie de Langue Française 2024 : 6 Abstracts publiés \newline dans le journal du congrès
\item Congrès National de Gériatrie et 4ème EUROMED Geriatrics Meeting – Tunis \newline 2023 : Poster présenté et Abstract publié dans le Middle East Journal of Age and \newline Ageing (Titre du poster : « COVID-19 transmission among elderly in Tunisia 2022-2023 »)
%\item[\textbf{[P.1]}] Inventor 1, Your Name, Inventor 3, et al. (Year). \href{https://patentoffice.gov/patent/XXXXXXXXX}{\textbf{Title of Patent}}. Patent Office, Patent No. XXXXXXXXX. Registration Date: Date, Grant Date: Date, Publication Date: Date.
%\item[\textbf{[J.1]}] Author 1, Your Name, Author 3, et al. (Year). \href{https://doi.org/XX.XXXX/XXXXX.XXXX.XXXXXXX}{\textbf{Title of Journal Article}}. \textit{Journal Name}, Vol. XX, Issue X, pp. XXX-XXX. DOI: XX.XXXX/XXXXX.XXXX.XXXXXXX
\end{itemize}
}
\vspace{-4mm}
\section{\textbf{EXPERIENCE EN LEADERSHIP}}
\vspace{-0.2mm}
\resumeSubHeadingListStart
\resumeProject
  {Participation à l’encadrement des externes en médecine }
  {Institut National Zouhaier Kallel de Nutrition}
  {2022 - 2023}
  {}
\resumeItemListStart
  \item Démonstrations pratiques des différentes compétences à savoir l’anamnèse, l’examen clinique,\newline la rédaction de dossiers médicaux, la discussion diagnostique et la conduite à tenir.
  \item Participer à la supervision des examens pratiques de type ECOS en tant qu’évaluatrice.
\resumeItemListEnd

\resumeProject
  {Responsable du comité culture }
  {Associa-Med Tunis ; Assoc. affiliée à la Fédération internationale des associations d’étudiants en médecine (IFMSA)}
  {2016 - 2017}

\resumeProject
  {Responsable logistique et recrutement des médecins dans la caravane de santé : Projet Doc'Time }
  {Associa-Med Tunis ; Assoc. affiliée à la Fédération internationale des associations d’étudiants en médecine (IFMSA)}
  {2015 - 2016}

\resumeSubHeadingListEnd
\vspace{-5mm}

\section{\textbf{ADHÉSION ET BÉNÉVOLAT}}
\vspace{-0.2mm}
\resumeSubHeadingListStart
\resumeProject
  {Bénévole au sein de l’Hôpital Jean Talon }
  {Santé Québec Nord-de-l’Île-de-Montréal – Universitaire}
  {2026}
  {}
\resumeItemListStart
  \item Visite amicale aux patients
\resumeItemListEnd

\resumeProject
  {Membre de l’association des diplômés internationaux en médecine de l’Alberta (AIMGA)}
  {}
  {2025}
 % {{}[\href{https://volunteer-org-a-link.com}{\textcolor{darkblue}{\faIcon{globe}}}]}

%\resumeProject
%  {Sensibilisation prevention et dépistage du SIDA / cancer du sein}
%  {AIMGA}
%  {2024}
%  {} %Ministere de la sante de Tunis
\resumeProject{Participation à des activités de sensibilisation et au dépistage du cancer du sein (Octobre rose)
}{}
  {2024}

\resumeProject{Participation à des activités de sensibilisation et au dépistage du Sida.
}{}
  {2024}

\resumeProject
  {Enquête épidémiologique sur la prévalence des troubles de l’humeur dans le gouvernorat de Mannouba}
  {en collaboration avec l'OMS}
  {2021}
  {}%OMS
\resumeItemListStart
  \item Participation en tant qu’enquêtrice
\resumeItemListEnd

\resumeProject
  {Membre actif de l’Associa-Med bureau de Tunis}
  {Faculté de Médecine de Tunis}
  {2013 - 2018}
  {}
\resumeItemListStart
  \item Accueil des nouveaux étudiants de première année pour améliorer leur intégration
  \item Implication dans des projets de sensibilisation, de promotion de la santé et d’actions \newline de terrain comme les caravanes de santé, en collaboration avec les étudiants et les professionnels de santé.
  \item Accompagnement de nourrissons abandonnés en milieu hospitalier en collaboration avec les intervenants concernés .
\resumeItemListEnd
\resumeSubHeadingListEnd
\vspace{-5mm}

%%%%%%%%%%%%%%%%%%%%%%%%%%%%%%%%%%%%%%%%%%%%%%%%%%%%%%%%%%%%%%%%%%%%%%%%%%%%%%%%%%%%%%%%%%%
\vspace{-0.4mm}
\section{\textbf{FORMATIONS}}
\vspace{-0.2mm}
\resumeSubHeadingListStart

\resumeProject
  {Advanced Life Support (ALS) }
  {European Resuscitation Council et Tunisian Resuscitation Council}
  {2025}

\resumeProject
  {Formation sur la méthodologie et la communication de l’évaluation rapide des risques de l’ECDC}
  {}
  {2023}

\vspace{-2mm}

\resumeProject
  {Masterclass sur Le bien-être au travail }
  {Faculté de médecine de Tunis}
  {2023}

\resumeProject
  {Formation certifiante « Wellbeing for Health Professional »}
  {University of Washington Global Health}
  {2022 - 2023}

\resumeProject
  {Masterclass « Le stigma dans le domaine de la santé – une formation pour comprendre et prévenir »}
  {Faculté de médecine de Tunis}
  {2022}

\resumeProject
  {Formation « Communication Skills and Facilitation » et « Project Management »}
  {}
  {2013}

\resumeSubHeadingListEnd

\vspace{-8mm}
\section{\textbf{Langues}}
\vspace{-0.4mm}
 \resumeHeadingSkillStart
  \resumeSubItem{Français: maîtrise professionnelle complète}
    {}
  \resumeSubItem{Anglais: courant}
    {}
  \resumeSubItem{Arabe: langue maternelle}
    {}
 \resumeHeadingSkillEnd

\vspace{-1mm}

\section{\textbf{COMPÉTENCES INFORMATIQUES}}
\vspace{-0.4mm}
 \resumeHeadingSkillStart
  \resumeSubItem{Dossier médical informatisé OACIS}
    {}
  \resumeSubItem{Microsoft office}
    {}
  \resumeSubItem{IBM SPSS Statistics}
    {}
  \resumeSubItem{Zotero}
    {}
 \resumeHeadingSkillEnd

\vspace{-1mm}

\section{\textbf{CENTRES D’INTÉRÊT}}
\vspace{-0.4mm}
 \resumeHeadingSkillStart
  \resumeSubItem{Handball / Randonnées / camping / Nature / Jeux de sociétés}
    {}
 \resumeHeadingSkillEnd


\end{document}